\begingroup
\renewcommand{\arraystretch}{1.45}
\fancyhead[L]{\small BESTIMMUNG DER ANZIEHUNG}
\Needspace{29\baselineskip}
\AttractionArticle{8}
Obwohl die Formeln 15, 16 zur Bestimmung der Koeffizienten $\gamma,\gamma',\gamma''$ ausreichen könnten, lassen sich doch noch elegantere angeben. Dazu werden wir die Gleichung [5] mit $aabb-GG$ multiplizieren, woraus sich nach einer kleinen Umformung ergibt
\[
\begin{aligned}
&\frac{aaAA(bb+G)}{aa+G}-AAG+\frac{bbBB(aa+G)}{bb+G}-BBG+\frac{aabbCC}{G}-CCG\\[6pt]
&\qquad{}=aabb-GG
\end{aligned}
\]
Aus der Beschaffenheit der kubischen Gleichung folgt aber
\[
\begin{aligned}
&\text{Summe der Wurzeln}\quad G-G'-G''=AA+BB+CC-aa-bb\\[8pt]
&\text{Produkt der Wurzeln}\quad GG'G''=aabbCC
\end{aligned}
\]
\Needspace{24\baselineskip}
Dadurch geht die vorangehende Gleichung in die folgende über:
\[
\begin{aligned}
&\frac{aaAA(bb+G)}{aa+G}+\frac{bbBB(aa+G)}{bb+G}+G'G''-G(G-G'-G''+aa+bb)\\[6pt]
&\qquad{}=aabb-GG
\end{aligned}
\]
die man auch so darstellen kann
\[
\begin{aligned}
&\frac{aaAA(bb+G)}{aa+G}+\frac{bbBB(aa+G)}{bb+G}-(aa+G)(bb+G)\\[6pt]
&\qquad{}+(G+G')(G+G'')=0
\end{aligned}
\]
\Needspace{30\baselineskip}
Hierdurch wird der Wert des Koeffizienten $\gamma$ aus der ersten Formel in [15] in den folgenden umgewandelt:
\[
\begin{aligned}
&\gamma=\surd{\frac{(aa+G)(bb+G)}{(G+G')(G+G'')}}\quad.\,.\,.\,.\,.\,.\,.\,.\quad[17]
\end{aligned}
\]
Durch eine ganz ähnliche Analysis findet man
\[
\begin{aligned}
&\gamma'=\surd{\frac{(aa-G')(bb-G')}{(G+G')(G''-G')}}\quad.\,.\,.\,.\,.\,.\,.\,.\quad[18]
\end{aligned}
\]
\[
\begin{aligned}
&\gamma''=\surd{\frac{(aa-G'')(bb-G'')}{(G+G'')(G'-G'')}}\quad.\,.\,.\,.\,.\,.\,.\quad[19]
\end{aligned}
\]
\SourceLeaf{354}
\par Nachdem die Koeffizienten $\gamma,\gamma',\gamma''$ gefunden sind, werden die übrigen $\alpha,\beta,\alpha',\beta',\alpha'',\beta''$ daraus durch die Formeln 3, 4, 6, 7, 9, 10 hergeleitet werden.

\Needspace{43\baselineskip}
\AttractionArticle{9}
Man sieht leicht, dass die Vorzeichen der Wurzelausdrücke, durch die wir $\gamma,\gamma',\gamma''$ bestimmt haben, nach Belieben gewählt werden können. Es lohnt sich aber zu untersuchen, wie das Vorzeichen der Größe $\varepsilon$ mit diesen Vorzeichen zusammenhängt. Dazu betrachten wir die dritte Gleichung in III, Art. 3.
\[
\begin{aligned}
&\varepsilon\gamma=\alpha'\beta''-\beta'\alpha''
\end{aligned}
\]
die durch die Formeln 6, 7, 9, 10 in diese übergeht:
\[
\begin{aligned}
&\varepsilon\gamma=\frac{abAB\gamma'\gamma''}{(aa-G')(bb-G'')}-\frac{abAB\gamma'\gamma''}{(aa-G'')(bb-G')}\\[8pt]
&=\frac{ab(aa-bb)AB(G''-G')\gamma'\gamma''}{(aa-G')(aa-G'')(bb-G')(bb-G'')}
\end{aligned}
\]
Aus der Betrachtung der Gleichung 13 leiten wir aber leicht ab
\[
\begin{aligned}
&(aa+G)(aa-G')(aa-G'')=aa(aa-bb)AA\\[8pt]
&(bb+G)(bb-G')(bb-G'')=-bb(aa-bb)BB
\end{aligned}
\]
\Needspace{22\baselineskip}
Dadurch wird die vorangehende Gleichung zu
\[
\begin{aligned}
&\varepsilon\gamma=\frac{(aa+G)(bb+G)(G'-G'')\gamma'\gamma''}{ab(aa-bb)AB}
\end{aligned}
\]
die, mit der Gleichung 17 verbunden,
\[
\begin{aligned}
&\gamma\gamma'\gamma''=\frac{\varepsilon ab(aa-bb)AB}{(G+G')(G+G'')(G'-G'')}
\end{aligned}
\]
liefert. \Needspace{33\baselineskip}
Hieraus ist ersichtlich, dass, wenn für $-G'$ die dem Betrag nach größere negative Wurzel der kubischen Gleichung gewählt ist und zugleich die Koeffizienten $\gamma,\gamma',\gamma''$ sämtlich positiv genommen sind, $\varepsilon$ dasselbe Vorzeichen erhält wie $AB$; und dass dasselbe geschieht, wenn bei allen vier Bedingungen oder bei zweien von ihnen die entgegengesetzte Wahl getroffen ist, das Gegenteil aber, wenn du eine oder drei Bedingungen umgekehrt hast. Außerdem sind noch die folgenden Beziehungen zu beachten, die sich leicht aus den vorangehenden herleiten lassen:
\SourceLeaf{355}
\[
\begin{aligned}
&\alpha\alpha'\alpha''=\frac{\varepsilon aabAAB}{(G+G')(G+G'')(G'-G'')}\\[8pt]
&\beta\beta'\beta''=-\frac{\varepsilon abbABB}{(G+G')(G+G'')(G'-G'')}\\[8pt]
&\alpha\beta=\frac{abAB}{(G+G')(G+G'')}\\[8pt]
&\alpha'\beta'=-\frac{abAB}{(G+G')(G'-G'')}\\[8pt]
&\alpha''\beta''=\frac{abAB}{(G+G'')(G'-G'')}
\end{aligned}
\]

\clearpage
\endgroup
